% OpenStax Calculus Volume 3 -- Section 6.5 Exercises: Divergence and Curl
% Exercises reproduced from OpenStax, licensed under CC BY 4.0.
% Source: https://openstax.org/books/calculus-volume-3/pages/6-5-divergence-and-curl
\documentclass{exercises}
\title{OpenStax Calculus Volume 3}
\subtitle{Section 6.5 Exercises: Divergence and Curl}
\sourceurl{https://openstax.org/books/calculus-volume-3/pages/6-5-divergence-and-curl}
\author{}
\date{}

\begin{document}
\maketitle


For the following exercises, determine whether the statement is true or false.

\begin{enumerate}[start=1]
\item If the coordinate functions of $\mathbf{F}:\mathbb{R}^{3}\to \mathbb{R}^{3}$ have continuous second partial derivatives, then $\operatorname{div}(\operatorname{curl}(\mathbf{F}))$ equals zero.
\item $\nabla \cdot(x\mathbf{i}+y\mathbf{j}+z\mathbf{k})=1$.
\item All vector fields of the form $\mathbf{F}(x,y,z)=f(x)\mathbf{i}+g(y)\mathbf{j}+h(z)\mathbf{k}$ are conservative.
\item If $\operatorname{curl}\mathbf{F}=0$, then $\mathbf{F}$ is conservative.
\item If $\mathbf{F}$ is a constant vector field, then $\operatorname{div}\mathbf{F}=0$.
\item If $\mathbf{F}$ is a constant vector field, then $\operatorname{curl}\mathbf{F}=0$.
\end{enumerate}

For the following exercises, find the curl of F.

\begin{enumerate}[resume]
\item $\mathbf{F}(x,y,z)=xy^{2}z^{4}\mathbf{i}+(2x^{2}y+z)\mathbf{j}+y^{3}z^{2}\mathbf{k}$
\item $\mathbf{F}(x,y,z)=x^{2}z\mathbf{i}+y^{2}x\mathbf{j}+(y+2z)\mathbf{k}$
\item $\mathbf{F}(x,y,z)=3xyz^{2}\mathbf{i}+y^{2}\sin z\mathbf{j}+xe^{2z}\mathbf{k}$
\item $\mathbf{F}(x,y,z)=x^{2}yz\mathbf{i}+xy^{2}z\mathbf{j}+xyz^{2}\mathbf{k}$
\item $\mathbf{F}(x,y,z)=(x\cos y)\mathbf{i}+xy^{2}\mathbf{j}$
\item $\mathbf{F}(x,y,z)=(x-y)\mathbf{i}+(y-z)\mathbf{j}+(z-x)\mathbf{k}$
\item $\mathbf{F}(x,y,z)=xyz\mathbf{i}+x^{2}y^{2}z^{2}\mathbf{j}+y^{2}z^{3}\mathbf{k}$
\item $\mathbf{F}(x,y,z)=xy\mathbf{i}+yz\mathbf{j}+xz\mathbf{k}$
\item $\mathbf{F}(x,y,z)=x^{2}\mathbf{i}+y^{2}\mathbf{j}+z^{2}\mathbf{k}$
\item $\mathbf{F}(x,y,z)=ax\mathbf{i}+by\mathbf{j}+c\mathbf{k}$ for constants $a$, $b$, $c$
\end{enumerate}

For the following exercises, find the divergence of F.

\begin{enumerate}[resume]
\item $\mathbf{F}(x,y,z)=x^{2}z\mathbf{i}+y^{2}x\mathbf{j}+(y+2z)\mathbf{k}$
\item $\mathbf{F}(x,y,z)=3xyz^{2}\mathbf{i}+y^{2}\sin z\mathbf{j}+xe^{2z}\mathbf{k}$
\item $\mathbf{F}(x,y)=(\sin x)\mathbf{i}+(\cos y)\mathbf{j}$
\item $\mathbf{F}(x,y,z)=x^{2}\mathbf{i}+y^{2}\mathbf{j}+z^{2}\mathbf{k}$
\item $\mathbf{F}(x,y,z)=(x-y)\mathbf{i}+(y-z)\mathbf{j}+(z-x)\mathbf{k}$
\item $\mathbf{F}(x,y)=\frac{x}{\sqrt{x^{2}+y^{2}}}\mathbf{i}+\frac{y}{\sqrt{x^{2}+y^{2}}}\mathbf{j}$
\item $\mathbf{F}(x,y)=x\mathbf{i}-y\mathbf{j}$
\item $\mathbf{F}(x,y,z)=ax\mathbf{i}+by\mathbf{j}+c\mathbf{k}$ for constants $a$, $b$, $c$
\item $\mathbf{F}(x,y,z)=xyz\mathbf{i}+x^{2}y^{2}z^{2}\mathbf{j}+y^{2}z^{3}\mathbf{k}$
\item $\mathbf{F}(x,y,z)=xy\mathbf{i}+yz\mathbf{j}+xz\mathbf{k}$
\end{enumerate}

For the following exercises, determine whether each of the given scalar functions is harmonic.

\begin{enumerate}[resume]
\item $u(x,y,z)=e^{-x}(\cos y-\sin y)$
\item $w(x,y,z)={(x^{2}+y^{2}+z^{2})}^{-1/2}$
\item If $\mathbf{F}(x,y,z)=2\mathbf{i}+2x\mathbf{j}+3y\mathbf{k}$ and $\mathbf{G}(x,y,z)=x\mathbf{i}-y\mathbf{j}+z\mathbf{k}$, find $\operatorname{curl}(\mathbf{F}\times \mathbf{G})$.
\item If $\mathbf{F}(x,y,z)=2\mathbf{i}+2x\mathbf{j}+3y\mathbf{k}$ and $\mathbf{G}(x,y,z)=x\mathbf{i}-y\mathbf{j}+z\mathbf{k}$, find $\operatorname{div}(\mathbf{F}\times \mathbf{G})$.
\item Find $\operatorname{div}\mathbf{F}$, given that $\mathbf{F}=\nabla f$, where $f(x,y,z)=xy^{3}z^{2}$.
\item Find the divergence of F for vector field $\mathbf{F}(x,y,z)=(y^{2}+z^{2})(x+y)\mathbf{i}+(z^{2}+x^{2})(y+z)\mathbf{j}+(x^{2}+y^{2})(z+x)\mathbf{k}$.
\item Find the divergence of F for vector field $\mathbf{F}(x,y,z)=f_{1}(y,z)\mathbf{i}+f_{2}(x,z)\mathbf{j}+f_{3}(x,y)\mathbf{k}$.
\end{enumerate}

For the following exercises, use $r=||\mathbf{r}||$ and $\mathbf{r}=\left\langle x,y,z \right\rangle$.

\begin{enumerate}[resume]
\item Find $\operatorname{curl}\mathbf{r}$.
\item Find $\operatorname{curl}\frac{\mathbf{r}}{r}$.
\item Find $\operatorname{curl}\frac{\mathbf{r}}{r^{3}}$.
\item Let $\mathbf{F}(x,y)=\frac{-y\mathbf{i}+x\mathbf{j}}{x^{2}+y^{2}}$, where $\mathbf{F}$ is defined on $\{(x,y)\in \mathbb{R}^{2}|(x,y)\ne(0,0)\}$. Find $\operatorname{curl}\mathbf{F}$.
\end{enumerate}

For the following exercises, use a computer algebra system to find the curl of the given vector fields.

\begin{enumerate}[resume]
\item \techrequired $\mathbf{F}(x,y,z)=\arctan(\frac{x}{y})\mathbf{i}+\ln \sqrt{x^{2}+y^{2}}\mathbf{j}+\mathbf{k}$
\item \techrequired $\mathbf{F}(x,y,z)=\sin(x-y)\mathbf{i}+\sin(y-z)\mathbf{j}+\sin(z-x)\mathbf{k}$
\end{enumerate}

For the following exercises, find the divergence of F at the given point.

\begin{enumerate}[resume]
\item $\mathbf{F}(x,y,z)=\mathbf{i}+\mathbf{j}+\mathbf{k}$ at $(2,-1,3)$
\item $\mathbf{F}(x,y,z)=xyz\mathbf{i}+y\mathbf{j}+x\mathbf{k}$ at $(1,2,3)$
\item $\mathbf{F}(x,y,z)=e^{-xy}\mathbf{i}+e^{xz}\mathbf{j}+e^{yz}\mathbf{k}$ at $(3,2,0)$
\item $\mathbf{F}(x,y,z)=xyz\mathbf{i}+y\mathbf{j}+z\mathbf{k}$ at $(1,2,1)$
\item $\mathbf{F}(x,y,z)=e^{x}\sin y\mathbf{i}-e^{x}\cos y\mathbf{j}$ at $(0,0,3)$
\end{enumerate}

For the following exercises, find the curl of F at the given point.

\begin{enumerate}[resume]
\item $\mathbf{F}(x,y,z)=\mathbf{i}+\mathbf{j}+\mathbf{k}$ at $(2,-1,3)$
\item $\mathbf{F}(x,y,z)=xyz\mathbf{i}+y\mathbf{j}+x\mathbf{k}$ at $(1,2,3)$
\item $\mathbf{F}(x,y,z)=e^{-xy}\mathbf{i}+e^{xz}\mathbf{j}+e^{yz}\mathbf{k}$ at $(3,2,0)$
\item $\mathbf{F}(x,y,z)=xyz\mathbf{i}+y\mathbf{j}+z\mathbf{k}$ at $(1,2,1)$
\item $\mathbf{F}(x,y,z)=e^{x}\sin y\mathbf{i}-e^{x}\cos y\mathbf{j}$ at $(0,0,3)$
\item Let $\mathbf{F}(x,y,z)=(3x^{2}y+az)\mathbf{i}+x^{3}\mathbf{j}+(3x+3z^{2})\mathbf{k}$. For what value of $a$ is $\mathbf{F}$ conservative?
\item Given vector field $\mathbf{F}(x,y)=\frac{1}{x^{2}+y^{2}}\left\langle-y,x \right\rangle$ on domain $D=\mathbb{R}^{2}-\{(0,0)\}=\{(x,y)\in \mathbb{R}^{2}|(x,y)\ne(0,0)\}$, is $\mathbf{F}$ conservative?
\item Given vector field $\mathbf{F}(x,y)=\frac{1}{x^{2}+y^{2}}\left\langle x,y \right\rangle$ on domain $D=\mathbb{R}^{2}-\{(0,0)\}=\{(x,y)\in \mathbb{R}^{2}|(x,y)\ne(0,0)\}$, is $\mathbf{F}$ conservative?
\item Find the work done by force field $\mathbf{F}(x,y)=e^{-y}\mathbf{i}-xe^{-y}\mathbf{j}$ in moving an object from $P(0,1)$ to $Q(2,0)$. Is the force field conservative?
\item Compute the divergence of $\mathbf{F}=(\sinh x)\mathbf{i}+(\cosh y)\mathbf{j}-xyz\mathbf{k}$.
\item Compute the curl of $\mathbf{F}=(\sinh x)\mathbf{i}+(\cosh y)\mathbf{j}-xyz\mathbf{k}$.
\end{enumerate}

For the following exercises, consider a rigid body that is rotating about the $x$-axis counterclockwise with constant angular velocity $\omega=\langle a,b,c\rangle$. If $P$ is a point in the body located at $\mathbf{r}=x\mathbf{i}+y\mathbf{j}+z\mathbf{k}$, the velocity at $P$ is given by vector field $\mathbf{F}=\omega \times \mathbf{r}$.


\textit{[Figure: A three dimensional diagram of an object rotating about the $x$-axis in a counterclockwise manner with constant angular velocity w = <a,b,c>. The object is roughly a sphere with pointed ends on the $x$-axis, which cuts it in half. An arrow r is drawn from (0,0,0) to P(x,y,z) and down from P(x,y,z) to the $x$-axis.]}

\begin{enumerate}[resume]
\item Express $\mathbf{F}$ in terms of $\mathbf{i}$, $\mathbf{j}$, and $\mathbf{k}$ vectors.
\item Find the divergence of $\mathbf{F}$.
\item Find the curl of $\mathbf{F}$.
\end{enumerate}

In the following exercises, suppose that $\nabla \cdot \mathbf{F}=0$ and $\nabla \cdot \mathbf{G}=0$.

\begin{enumerate}[resume]
\item Does $\mathbf{F}+\mathbf{G}$ necessarily have zero divergence?
\item Does $\mathbf{F}\times \mathbf{G}$ necessarily have zero divergence?
\end{enumerate}

In the following exercises, suppose a solid object in $\mathbb{R}^{3}$ has a temperature distribution given by $T(x,y,z)=100e^{-x^{2}-y^{2}-z^{2}}$. The heat flow vector field in the object is $\mathbf{F}=-k\nabla T$, where $k>0$ is a property of the material. The heat flow vector points in the direction opposite to that of the gradient, which is the direction of greatest temperature decrease. The divergence of the heat flow vector is $\nabla \cdot \mathbf{F}=-k\nabla \cdot \nabla T=-k\nabla^{2}T$.

\begin{enumerate}[resume]
\item Compute the heat flow vector field.
\item Compute the divergence.
\item \techrequired Consider rotational velocity field $\mathbf{v}=\langle0,10z,-10y\rangle$. If a paddlewheel is placed in plane $x+y+z=1$ with its axis normal to this plane, using a computer algebra system, calculate how fast the paddlewheel spins in revolutions per unit of time. \\ \textit{[Figure: A three dimensional diagram of a rotational velocity field. The arrows are showing a rotation in a clockwise manner. A paddlewheel is shown in plane x + y + z = 1 with n extended out perpendicular to the plane.]}
\end{enumerate}

\end{document}
